\begin{theorem} \label{thm:t1} 
Sea $A \subseteq \mathbb{R}^n$ un conjunto abierto y arcoconexo (ver la Definición~\ref{def:conjuntoconexo_simp}) y sea $f: A \subseteq \mathbb{R}^n \to \mathbb{R}$ un campo de clase $\mathcal{C}^1(A)$. Además, sea $\boldsymbol{\sigma}: [a,b] \to A \subseteq \mathbb{R}^n$ una trayectoria de clase $\mathcal{C}^{1}$ a trozos. Entonces: 
\[ \int_{\boldsymbol{\sigma}} \nabla f \cdot d\mathbf{s} = f(\boldsymbol{\sigma}(b)) - f(\boldsymbol{\sigma}(a)) \] 
\end{theorem}

\begin{proof}
Demostraremos el caso en el cual $\boldsymbol{\sigma}$ es una trayectoria de clase $\mathcal{C}^{1}$ y se deja al lector probar el caso más general. 

Definamos la función escalar $g(t) = (f \circ \boldsymbol{\sigma})(t)$ con $t \in [a,b]$. Sabemos que $g$ resulta derivable y, por la regla de la cadena, obtenemos: 
\begin{equation*} 
g'(t) = (f \circ \boldsymbol{\sigma})'(t) = \nabla f(\boldsymbol{\sigma}(t)) \cdot \boldsymbol{\sigma}'(t) 
\end{equation*} 

Por el Teorema Fundamental del Cálculo de una variable, sabemos que: 
\begin{equation*} 
\int_{a}^{b} g'(t) \, dt = g(b) - g(a) = f(\boldsymbol{\sigma}(b)) - f(\boldsymbol{\sigma}(a)) 
\end{equation*} 

Por consiguiente:
\begin{align*} 
\int_{\boldsymbol{\sigma}} \nabla f \cdot d\mathbf{s} &= \int_{a}^{b} \nabla f(\boldsymbol{\sigma}(t)) \cdot \boldsymbol{\sigma}'(t) \, dt \\ 
&= \int_{a}^{b} g'(t) \, dt = g(b) - g(a) \\ 
&= f(\boldsymbol{\sigma}(b)) - f(\boldsymbol{\sigma}(a))   \qedhere
\end{align*} 
\end{proof}



\begin{tarea} Para la trayectoria  $\boldsymbol{\sigma}:[0,2]\to A\subseteq\Rn{3}\; : \; \sigma(t) = (t^{3}+1, t^{2}, t^{3}),$  calcular  $$ \int_{\sigma}  (y, x, 0)  \cdot d\mathbf{s} $$
\end{tarea}


\begin{corollary}\label{thm:c1} 
Sean $A \subseteq \mathbb{R}^n$ un conjunto abierto y arcoconexo y $f: A \subseteq \mathbb{R}^n \to \mathbb{R}$ un campo de clase $\mathcal{C}^1(A)$. Si $\boldsymbol{\sigma}_1: [a_1,b_1] \to A$ y $\boldsymbol{\sigma}_2: [a_2,b_2] \to A$ son dos trayectorias de clase $\mathcal{C}^{1}$ a trozos tales que $\boldsymbol{\sigma}_1(a_1) = \boldsymbol{\sigma}_2(a_2) \ \land \ \boldsymbol{\sigma}_1(b_1) = \boldsymbol{\sigma}_2(b_2)$, entonces: 
\[ \int_{\boldsymbol{\sigma}_1} \nabla f \cdot d\mathbf{s} = \int_{\boldsymbol{\sigma}_2} \nabla f \cdot d\mathbf{s} \] 
Es decir, la integral $\int_{\boldsymbol{\sigma}} \nabla f \cdot d\mathbf{s}$ sólo depende de los puntos inicial y final de la trayectoria. 
\end{corollary}

\begin{proof}
Demostraremos el caso en el cual $\boldsymbol{\sigma}_1$ y $\boldsymbol{\sigma}_2$ son trayectorias de clase $\mathcal{C}^{1}$ y se deja al lector probar el caso más general. 

Usando el Teorema~\ref{thm:t1}, se tiene que:
\begin{align*} 
\int_{\boldsymbol{\sigma}_1} \nabla f \cdot d\mathbf{s} &= f(\boldsymbol{\sigma}_1(b_1)) - f(\boldsymbol{\sigma}_1(a_1)) \\ 
\int_{\boldsymbol{\sigma}_2} \nabla f \cdot d\mathbf{s} &= f(\boldsymbol{\sigma}_2(b_2)) - f(\boldsymbol{\sigma}_2(a_2)) 
\end{align*} 

Como por hipótesis $\boldsymbol{\sigma}_1(a_1) = \boldsymbol{\sigma}_2(a_2) \ \land \ \boldsymbol{\sigma}_1(b_1) = \boldsymbol{\sigma}_2(b_2)$, por la buena definición de la función $f$ se cumple que: 
\begin{equation*} 
f(\boldsymbol{\sigma}_1(a_1)) = f(\boldsymbol{\sigma}_2(a_2)) \ \land \ f(\boldsymbol{\sigma}_1(b_1)) = f(\boldsymbol{\sigma}_2(b_2)) 
\end{equation*} 

Luego, restando miembro a miembro obtenemos:
\begin{equation*} 
\int_{\boldsymbol{\sigma}_1} \nabla f \cdot d\mathbf{s} - \int_{\boldsymbol{\sigma}_2} \nabla f \cdot d\mathbf{s} = 0,
\end{equation*} 
lo que demuestra el resultado buscado. 
\end{proof}






